\section{Conclusion and Outlook}
\label{section:Conclusion}





The goal was to find a device design and fabrication method that would result in optimal superconductance at the tips of the leads of the SEMBJ. The most important parameter determining the quality of the superconductance is the size of the induced energy gap at the tips of the leads. The reason that this parameter is important is that it determines the voltage bias range in which the device operates in the superconducting regime; for bias voltages outside the range determined by this energy gap, transport happens normally and the reasons for using an SEMBJ over a regular EMBJ are lost. 

The space of device designs and fabrication methods was constrained to only include bowtie geometry constrictions fabricated using shadow mask double angle evaporation. The only superconductor that was evaluated was niobium. The parameters that were varied during the optimization of the design were the thicknesses of the bridge layer and the superconducting layer, and the length of the bridge. The devices were electromigrated using feedback controlled electromigration, both at room temperature and at low temperature. The devices were characterized both before and after electromigration using low temperature transport measurements and SEM imaging. The results of these characterizations were used to inform the subsequent designs during each design iteration. 

This chapter summarizes the main findings of this thesis, evaluates their significance, discusses their limitations and suggests possible avenues for future research. 


\subsection{Summary of Main Findings}

There are several important conclusions that can be drawn from the results in this thesis. The main conclusions are listed here. 

\begin{itemize}

\item It is possible to create proximity induced superconducting electromigrated breakjunctions using double angle shadow mask evaporation and feedback controlled electromigration in a bowtie geometry.

\item It is possible to proximitize palladium and gold films using this method.

\item Our experiments resulted in a proximity-induced energy gap of approximately 0.28 meV at the tips of the leads of the SEMBJ. This gap size is comparable to values reported in previous studies, indicating that our device design and fabrication method are effective for achieving high-quality superconductivity in this type of junction. The ability to simply and reliably produce proximity-induced gaps of this magnitude is an important step toward realizing practical applications of these devices and for future research.

\item Good control over the location of the electromigration-induced break in the device is necessary to ensure a large energy gap at the ends of both leads.

\item The results of this study demonstrate that bowtie constricted gold films with a thickness of 15 nm can be successfully electromigrated with a high degree of control. These findings suggest that it may be feasible to increase the thickness of the gold films while maintaining sufficient control over the electromigration process.

\item The results of this study suggest that a 50 nm thick niobium layer is capable of inducing superconductivity in a 30 nm thick gold layer, resulting in an energy gap within one percent of that of pure niobium. These findings imply that it may be possible to further increase the thickness of the gold while maintaining a sufficient bilayer energy gap.

\item The devices show an asymmetric critical current upon applying a magnetic field in the perpendicular in-plane direction. Our device design offers a simple lithographic platform for making Josephson-diode based devices.

\end{itemize}





\subsection{Suggested Future Research}

As mentioned, the resubstitution of palladium for the bridge material in the device design was not researched due to time constraints. Testing devices of this type would be an obvious next step in continuing this research.

Strong correlations were found between the $I_CR_N$ products of the devices and their normal state resistances. Having knowledge and control of such correlations could be a valuable prognostic tool, that could be used to pre-select especially promising devices based on simple room temperature measurements. Exploring this correlation further would therefore be of great interest.

The fabrication and measurements of the kinds of devices studied in this thesis is a slow process, which necessarily limits the number of design parameters that can be studied. Continued exploration of the possible device designs is necessary for further progress. An interesting design variation that could be explored in future work is an SEMBJ where the normal metal is sandwiched in between two superconducting layers in an SNS sandwich, instead of only being proximitized from one side in an SN bilayer design as was studied in this thesis. This design could be realized using the same shadow mask that was used during this thesis with no changes or only minimal changes, and seems intuitively better suited for inducing superconductivity than an SN bilayer design.

In this thesis we showed that our devices possess an asymmetric critical current upon applying a magnetic field in the perpendicular in-plane direction. Future research can build on this finding by exploring variations on the design. 

\subsection{Concluding Remarks}

The goal of this thesis, to create an SEMBJ, was reached. The fabrication method that was used throughout this thesis stands out due to its simplicity and swiftness; it contains only a single e-beam step, and many dozens of devices can be made in only a few hours. This is a significant advantage that allows for rapid device design changes and cheaper and faster experiments. Even when the design is changed to incorporate a local back gate, it remains simple compared to many alternative fabrication methods. The induced gap at the leads of the SEMBJ, which reached 280 $\mu$eV in our final design, is of the same order or larger than many devices found in the literature, indicating that our method is effective for achieving high-quality superconductivity in this type of junction. Additionally, the devices were found to have a magnetic field dependent directional asymmetry in their critical current at each stage of breaking. This phenomenon is currently not well understood, and the findings in this thesis may help guide future research.

